\begin{remark}[The R7 hypotheses used in this Part]
\label{c:r7-audit}
\status{Used and rechecked, with its internal-Hom counterexample confirmed}
We checked the proof in the supplied reading, \S\S2.1--2.3, against 12:150--158. The conditional construction there is not a construction of a faithful realization on all Chow motives: it first quotients by the chosen homological equivalence, assumes its equality with numerical equivalence and the stated positivity, and then takes the semisimple envelope. Our Chow assertions above do not use this conditional passage. In particular it cannot be used to kill the homologically zero K3 determinant.

The sign correction is forced by the cross-product formula: interchanging a class of degree $a$ with one of degree $b$ multiplies it by $(-1)^{ab}$. Thus the original realization is symmetric into graded vector spaces with that commutativity constraint. This is the convention used in our Frobenius equivalence. A level filtration by thick subcategories yields a direct-sum level decomposition only when the ambient category is semisimple; then we group the simple constituents according to their least allowed twist. Thickness alone gives no splitting of extensions. The geometric level definition of C.4 therefore does not assert a canonical splitting or identify itself with coniveau.

The internal-Hom error of 12:156 already occurs in the semisimple effective Tate category with objects finite sums of $\Q(-a)$, $a\geq0$. For every effective $Z$,
\[
 \Hom(Z\otimes\Q(-1),\Q)=0.
\]
Hence the effective internal Hom from $\Q(-1)$ to $\Q$ is the zero object. After adjoining all Tate twists the corresponding internal Hom is $\Q(1)$, which is nonzero. Thus effective internal Homs need not be preserved by Tate inversion. In a rigid tensor category the valid identity is
\[
 \underline{\Hom}(M,N(-i))
 =M^\vee\otimes N(-i)
 =M^*\otimes N,
 \qquad M^*=M^\vee(-i).
\]
Evaluation and coevaluation give these isomorphisms and their inverse maps. Merely postulating the effective object $M^*$ without the duality identities does not prove them. These checks confirm precisely the restrictions of R7 needed here.
\end{remark}
